\SourceLeaf{385}%
\textsc{von} \textsc{Zach} in 1809.

The following section contains several interesting observations and investigations connected with the state of the atmosphere, namely simultaneous meteorological observations at the hospice on Mont Cenis and in Milan, a barometric determination of the heights of the former station and Mont Colombier, a trigonometric determination of the heights of Mont-blanc (4802,7 metres) and Monte Rosa (4619,6 metres), and finally investigations of terrestrial refraction. The latter is determined from the angles of elevation of three mountains, Rocca Melone, Monte Viso and Monte Rosa, observed from three points (Milan, Turin, Mondovi). The absolute heights of the three stations are taken as given, and the heights of the observed points are eliminated. This is a procedure that can provide little certainty and that, on closer examination, also yielded little agreement. We would therefore attach little weight to the final result for the ratio of the Earth's curvature to the \textit{total} refraction (5,28 to 1). Calculation from six pairs of reciprocal zenith distances between principal triangulation points gives us a much smaller ratio on average, namely 1 to 0,1235, in very close agreement with the results found in the Hanoverian and Livonian surveys of arcs.

The tenth section deals with a much-discussed problem. Given the latitude of one endpoint of a given side of a triangle and the azimuth of that side at that endpoint, find the same quantities for the other endpoint and the difference in longitude on the elliptical spheroid. The derivation merely transforms \textsc{Legendre}'s formulas to introduce the true latitude in place of the so-called reduced latitude. The formulas are accurate through quantities of the third order, insofar as the flattening and the side of the triangle (taking the Earth's radius as the unit) are regarded as quantities of the first order. In applying the formulas to the present measurements, the quantities of the third order were omitted on the grounds that this is sufficiently accurate in practice. This can be conceded only insofar as the results are to be used merely for comparison with the astronomical determinations. For that purpose there is indeed no need to calculate the former with much greater precision than the latter permit. But if one takes a different point of view and calculates the geodetic results as accurately as
